\section{Block Encoding with Unbounded Spectra}
\label{sec:block-encoding}

The single-qubit theory of Section~\ref{sec:main} transforms a
scalar signal $r$ into an unbounded rational function.  To
apply this machinery to a multi-qubit Hamiltonian
$H = \sum_j \lambda_j\ket{j}\bra{j}$, one needs a way to
feed each eigenvalue $\lambda_j$ into the QSP circuit
simultaneously---the quantum-computing analogue of ``applying a
function to a matrix.''  In the standard QSVT
framework~\cite{GSLW2019}, this is accomplished by block
encoding, which embeds $H/\alpha$ (with $\alpha \geq \|H\|$)
into a larger unitary.  The normalization factor~$\alpha$
is the price of boundedness: it inflates the circuit depth and
forces the target polynomial to approximate $f(\alpha\,\cdot\,)$
on $[-1,1]$.

Stereographic encoding offers an alternative that avoids this
normalization entirely, at the cost of requiring knowledge of
the eigenbasis---a strong assumption that restricts the approach
to structured Hamiltonians.  This section defines the
stereographic block encoding, shows how the full QSP protocol
lifts to a multi-qubit eigenvalue transformation, and
demonstrates the construction through explicit circuits of
increasing complexity.

\subsection{Standard and Stereographic Block Encoding}
\label{subsec:block-encoding-defs}

To fix notation and set up the comparison, we begin with the
standard framework.

\begin{definition}[Standard Block Encoding]
	An $(\alpha, m, \epsilon)$-block encoding of a Hermitian operator $H$ acting on $n$ qubits is a unitary $U$ acting on $n + m$ qubits such that
	\begin{equation}
		\|H - \alpha(\bra{0}^{\otimes m} \otimes I)\, U\, (\ket{0}^{\otimes m} \otimes I)\| \leq \epsilon,
	\end{equation}
	where $\alpha \geq \|H\|$.
\end{definition}

To apply QSVT and compute $f(H)$, one constructs a degree-$d$
polynomial $P$ satisfying $|P(\lambda/\alpha)| \leq 1$ for all
eigenvalues $\lambda$, then finds QSP phases for $P$.  This
approach incurs two costs: the subnormalization factor
$\alpha \geq \|H\|$ increases circuit depth by a factor
of~$\alpha$, and the polynomial $P$ must approximate
$f(\alpha\,\cdot\,)$ on $[-1,1]$, which may require high degree
if $f$ has sharp features or singularities.

The stereographic alternative removes both costs by encoding
each eigenvalue directly, without normalization.

\begin{definition}[Stereographic Block Encoding]
	\label{def:stereo-block}
	Let $H$ be a Hermitian operator with spectral decomposition $H = \sum_j \lambda_j \ket{j}\bra{j}$, where $\lambda_j \geq 0$ (a positivity assumption that can be relaxed by a shift). A stereographic block encoding of $H$ is a unitary $U$ acting on $n + 1$ qubits such that
	\begin{equation}
		U = \sum_j U_{\lambda_j} \otimes \ket{j}\bra{j},
		\label{eq:stereo-block}
	\end{equation}
	where $U_{\lambda_j}$ is the single-qubit encoding unitary~\eqref{eq:encoding-unitary} with $r = \lambda_j$ and $\theta = 0$.
\end{definition}

Explicitly, in the eigenspace of eigenvalue $\lambda_j$:
\begin{equation}
	U_{\lambda_j} = \frac{1}{\sqrt{1+\lambda_j^2}}\begin{pmatrix} \lambda_j & 1 \\ 1 & -\lambda_j \end{pmatrix},
\end{equation}
which maps $\ket{0}\ket{j} \mapsto \ket{\lambda_j}\ket{j}$ using
a single ancilla qubit.  Each eigenvalue, no matter how large,
is encoded into a point on the Bloch sphere via stereographic
projection---the same compression
$\lambda_j \mapsto \tilde{\lambda}_j =
	\lambda_j/\sqrt{1+\lambda_j^2} \in [0,1)$ that drives the
single-qubit theory.

\begin{remark}[Caveats and scope]
	\label{rem:caveats}
	Constructing~\eqref{eq:stereo-block} requires the ability to apply $U_{\lambda_j}$ conditioned on the eigenspace $\ket{j}$ of $H$---in other words, it assumes access to the eigenbasis. This is a strong assumption that is \emph{not} met for generic Hamiltonians. We do not claim that stereographic block encoding is universally applicable; rather, it applies naturally to structured operators whose eigenbasis is known or efficiently computable.
\end{remark}

\subsection{The Stereographic QSVT Protocol}
\label{subsec:stereo-qsvt}

With the block encoding in place, the QSP protocol from
Section~\ref{subsec:qsp-phases} lifts to a multi-qubit
eigenvalue transformation.  For a Hamiltonian with known
diagonalization $H = V^\dagger D V$, the protocol interleaves
the signal operator with phase rotations within the eigenbasis:
\begin{equation}
	\begin{quantikz}[column sep=0.2cm, row sep=0.4cm]
		\lstick{\footnotesize$\ket{0}_a$} & \gate{\phi_d} & \gate{W_j} & \ \cdots\ \qw & \gate{W_j} & \gate{\phi_0} & \meter{} \\
		\lstick{\footnotesize$\ket{\psi}_s$} & \gate[2]{V} & \ctrl{-1}\vqw{1} & \ \cdots\ \qw & \ctrl{-1}\vqw{1} & \gate[2]{V^\dagger} & \qw \\
		& & \ctrl{-1} & \ \cdots\ \qw & \ctrl{-1} & & \qw
	\end{quantikz}
	\label{eq:full-protocol-circuit}
\end{equation}
Each $\phi_k$ box denotes $R_z(\phi_k) = \exp[i\phi_k\PauliZ]$,
and $W_j \equiv W(\tilde{\lambda}_j)$ is the stereographic
signal unitary on the ancilla conditioned on eigenstate
$\ket{j}$.  Because the signal operator acts independently in
each eigenspace, a degree-$d$ QSP sequence with phases
$\bm{\phi}$ produces
\begin{equation}
	W_{\bm{\phi}}(\lambda_j)\ket{0}\ket{j} = \left(P(\tilde{\lambda}_j)\ket{0} + Q(\tilde{\lambda}_j)\ket{1}\right)\ket{j},
\end{equation}
where $\tilde{\lambda}_j = \lambda_j/\sqrt{1+\lambda_j^2}$.
Decoding the ancilla via Pauli measurements yields the rational
eigenvalue transformation
\begin{equation}
	f(\lambda_j) = \frac{P(\tilde{\lambda}_j)}{Q(\tilde{\lambda}_j)}.
	\label{eq:eigenvalue-transform}
\end{equation}
The result is defined for all $\lambda_j \geq 0$ (including
$\lambda_j \gg 1$), can diverge wherever
$Q(\tilde{\lambda}_j) = 0$, and is a rational function of
$\lambda_j$ of total degree $\leq d$.  In contrast to standard
QSVT, no normalization factor~$\alpha$ enters the circuit: the
QSP degree~$d$ controls the approximation quality, and the
circuit depth scales as $\mathcal{O}(d)$ regardless of $\|H\|$.

Even without applying phase factors, the block encoding itself
already extracts useful spectral information.

\begin{proposition}[Extraction of Eigenvalues]
	\label{prop:extraction}
	Measuring the ancilla qubit in the Pauli basis after applying the stereographic block encoding $U$ to $\ket{0}_a\ket{j}$ yields the Bloch vector components
	\begin{equation}
		\langle\PauliZ\rangle = \frac{\lambda_j^2 - 1}{\lambda_j^2 + 1},
	\end{equation}
	from which the eigenvalue is recovered as $\lambda_j = \sqrt{(1+\langle\PauliZ\rangle)/(1-\langle\PauliZ\rangle)}$, without requiring knowledge of $\|H\|$.
\end{proposition}

This is the multi-qubit analogue of the stereographic decoding
formula~\eqref{eq:decoding}: the Bloch sphere coordinates of the
ancilla encode each eigenvalue independently, and no global
normalization constant is needed to extract them.

\subsection{Circuit Constructions}
\label{subsec:concrete-block}

The abstract definition~\eqref{eq:stereo-block} becomes concrete
once one specifies how to implement the eigenspace-conditioned
unitaries $U_{\lambda_j}$.  We present three examples of
increasing complexity: a diagonal Hamiltonian (where the
eigenbasis is the computational basis), a Pauli-$Z$ sum (where
algebraic structure reduces the gate count), and the two-qubit
Heisenberg model (where a non-trivial change of basis is
required).

\subsubsection{Diagonal Hamiltonians}
\label{subsubsec:diagonal}

Let $H$ be a diagonal Hamiltonian on $n$ qubits with
eigenvalues $\lambda_j \geq 0$ (which may be arbitrarily
large) in the computational basis.  This is the simplest case
because no diagonalization step is needed: the encoding
unitary acts directly on the ancilla, conditioned on the
computational-basis state $\ket{j}$ of the system register.

For each eigenvalue, the encoding state $\ket{\lambda_j}$ is
prepared from $\ket{0}$ by a $Y$-rotation:
\begin{equation}
	R_y(2\alpha_j)\ket{0} = U_{\lambda_j}\ket{0} = \ket{\lambda_j}, \qquad \alpha_j = \arctan(1/\lambda_j),
	\label{eq:diagonal-ry}
\end{equation}
since $\cos\alpha_j = \lambda_j/\sqrt{1+\lambda_j^2}$ and
$\sin\alpha_j = 1/\sqrt{1+\lambda_j^2}$.  The full encoding
unitary satisfies
$U_{\lambda_j} = R_y(2\alpha_j)\,\sigma_z$ (note
$\det U_{\lambda_j} = -1$), and the stereographic signal
operator~\eqref{eq:signal-operator} in each eigenspace becomes
$S_{\lambda_j} = U_{\lambda_j}\sigma_z = R_y(2\alpha_j)$,
a pure rotation.

The block encoding~\eqref{eq:stereo-block} is therefore a
uniformly controlled rotation: for each $j$, apply
$R_y(2\alpha_j)$ on the ancilla conditioned on the system
register being in state $\ket{j}$.  For a 2-qubit system
($N=4$ eigenvalues), the circuit is:
\begin{equation}
	\begin{quantikz}[column sep=0.5cm, row sep=0.5cm]
		\lstick{$\ket{0}_a$} & \gate[style={inner sep=4pt}]{R_y(2\alpha_j)} & \qw \\
		\lstick{$\ket{j_1}$} & \ctrl{-1}\vqw{1} & \qw \\
		\lstick{$\ket{j_0}$} & \ctrl{-1} & \qw
	\end{quantikz}
	\label{eq:diagonal-circuit}
\end{equation}
where the double control line indicates a uniformly controlled
rotation: the angle
$\alpha_j = \arctan(1/\lambda_j)$ is selected by the
computational-basis state $\ket{j} = \ket{j_1 j_0}$.
Standard decomposition techniques~\cite{NielsenChuang2010}
yield $\mathcal{O}(N)$ CNOT gates and $\mathcal{O}(N)$
single-qubit rotations, matching the gate complexity of
standard diagonal block encoding---but with a single ancilla
qubit rather than $\lceil\log_2\alpha\rceil + 1$, and with
rotation angles that depend only on individual eigenvalues,
not on $\max_j \lambda_j$.  Adding a new eigenvalue
$\lambda_{N}$ to the spectrum does not change the encoding of
any existing eigenvalue.

\subsubsection{Pauli-$Z$ Sum Hamiltonians}
\label{subsubsec:pauli-z}

When additional algebraic structure is present, the gate count
can be reduced below $\mathcal{O}(N)$.  Consider a Hamiltonian
that is a sum of tensor products of Pauli-$Z$ operators and
identity:
\begin{equation}
	H = \sum_{k=1}^{m} c_k\, Z_{i_k},
	\label{eq:pauli-z-ham}
\end{equation}
where $Z_{i_k}$ denotes $Z$ on qubit $i_k$ (and identity
elsewhere) and $c_k > 0$.  Such Hamiltonians include classical
Ising models and arise as the diagonal part of many physically
relevant systems.  Since each $Z_{i_k}$ has eigenvalues
$\pm 1$, the eigenvalues of $H$ are
$\lambda_{\mathbf{s}} = c_1 s_1 + \cdots + c_m s_m$ with
$s_k \in \{+1, -1\}$, and the Hamiltonian is already diagonal
in the computational basis, so the diagonal construction
applies directly.

However, the rotation angle
$\alpha_\mathbf{s} = \arctan(1/\lambda_\mathbf{s})$ depends on
$\lambda_\mathbf{s}$ only through the partial sums
$\sum c_k s_k$, which are determined by $m$ bits rather than
$n$ qubits.  This allows one to compute the partial sums into
$m$ ancilla qubits via a reversible adder, apply a
controlled-$R_y$ conditioned on the result, and uncompute.  The
gate complexity is $\mathcal{O}(\mathrm{poly}(m))$ rather than
$\mathcal{O}(2^n)$, making this efficient when $m \ll n$.

\begin{remark}[Positivity]
	Equation~\eqref{eq:pauli-z-ham} with $c_k > 0$ has eigenvalues in $[-\sum c_k, +\sum c_k]$, which includes negative values. To apply the stereographic encoding (which assumes $\lambda_j \geq 0$), shift by $\lambda_{\mathrm{shift}} = \sum_k c_k$ so that $H' = H + \lambda_{\mathrm{shift}} I \geq 0$. The decoded function then computes $f(\lambda + \lambda_{\mathrm{shift}})$, which must be accounted for in the target function design.
\end{remark}

\subsubsection{Two-Qubit Heisenberg Model}
\label{subsubsec:heisenberg}

The first two examples have computational-basis eigenstates. As a
non-trivial example requiring a genuine change of basis, consider
the isotropic Heisenberg model with a longitudinal field on two
qubits:
\begin{equation}
	H = J(X \otimes X + Y \otimes Y + Z \otimes Z) + h(Z \otimes I),
	\label{eq:heisenberg}
\end{equation}
where $J$ is the exchange coupling and $h$ is the field strength.
The Hamiltonian is block-diagonal in the standard basis:
$\ket{00}$ and $\ket{11}$ are eigenstates with eigenvalues
$\lambda_1 = J + h$ and $\lambda_4 = J - h$, while the
$\{\ket{01}, \ket{10}\}$ subspace has the $2\times 2$ block
\begin{equation}
	\begin{pmatrix} -J+h & 2J \\ 2J & -J-h \end{pmatrix}
\end{equation}
with eigenvalues
\begin{equation}
	\lambda_{2,3} = -J \pm \sqrt{h^2 + 4J^2}.
	\label{eq:heisenberg-spectrum}
\end{equation}
For $J = 1$, $h = 0.5$, the eigenvalues are
$\lambda \approx \{-3.06, 0.50, 1.06, 1.50\}$.  When $h = 0$,
the eigenstates in the $\{\ket{01},\ket{10}\}$ subspace reduce
to the Bell states $(\ket{01}\pm\ket{10})/\sqrt{2}$ with
eigenvalues $J$ and $-3J$.

The change-of-basis unitary $V$ that diagonalizes $H$ (i.e.,
$H = V^\dagger D V$ with
$D = \mathrm{diag}(\lambda_1, \ldots, \lambda_4)$) acts
non-trivially only on the $\{\ket{01}, \ket{10}\}$ subspace as
a Givens rotation, decomposable into $\mathcal{O}(1)$ elementary
gates.  The full stereographic block encoding on 3 qubits
(2 system + 1 ancilla) is then
\begin{equation}
	U = (V^\dagger \otimes I_a) \cdot U_{\mathrm{diag}} \cdot (V \otimes I_a),
	\label{eq:heisenberg-stereo}
\end{equation}
where $U_{\mathrm{diag}}$ is the diagonal stereographic encoding
from the first example, giving the circuit:
\begin{equation}
	\begin{quantikz}[column sep=0.4cm, row sep=0.5cm]
		\lstick{$\ket{0}_a$} & \qw & \gate[style={inner sep=4pt}]{R_y(2\alpha_j)} & \qw & \qw \\
		\lstick{$\ket{\psi_1}$} & \gate[2]{V} & \ctrl{-1}\vqw{1} & \gate[2]{V^\dagger} & \qw \\
		\lstick{$\ket{\psi_0}$} & \qw & \ctrl{-1} & \qw & \qw
	\end{quantikz}
	\label{eq:heisenberg-circuit}
\end{equation}
The three steps---diagonalize ($V$), encode
($U_{\mathrm{diag}}$), un-diagonalize ($V^\dagger$)---each
require $\mathcal{O}(1)$ gates, giving a total gate count that
is constant, independent of $\|H\|$, and uses a single ancilla
qubit.

\begin{remark}[Positivity shift for Heisenberg]
	For $J > 0$, the smallest eigenvalue $\lambda_{\min} = -J - \sqrt{h^2+4J^2}$ is negative. Shifting by $\lambda_{\mathrm{shift}} = |\lambda_{\min}|$ makes all eigenvalues non-negative, with the shifted minimum at zero. This zero eigenvalue maps to the south pole of the Bloch sphere ($\ket{1}$) under stereographic encoding and requires special handling for functions like $1/\lambda$ (e.g., regularization by a small shift $\delta > 0$).
\end{remark}




